\section{Thiết kế kiến trúc tổng thể}\label{sec:arch}

\subsection{Kiến trúc phân lớp}

Hotelia được thiết kế theo kiến trúc khách -- chủ nhiều lớp, với một máy chủ API dạng \emph{khối đơn thể phân module} (modular monolith). \Cref{fig:arch} mô tả bốn lớp: lớp trình bày gồm ứng dụng iOS và bảng điều khiển web; lớp API gồm máy chủ NestJS; lớp dữ liệu là MongoDB Atlas; và các dịch vụ bên ngoài gồm payOS, nhà cung cấp danh tính và Open-Meteo. Mô hình LightGBM không phải một dịch vụ riêng mà là một tệp JSON được nạp vào bộ nhớ của tiến trình NestJS khi khởi động.

\begin{figure}[H]
\centering
\diagram{arch-overall}
\caption{Kiến trúc tổng thể của Hotelia}\label{fig:arch}
\end{figure}

\Cref{tab:arch-decisions} tóm tắt các quyết định kiến trúc chính cùng lý do và đánh đổi. Quyết định quan trọng nhất là không tách vi dịch vụ. Với một nhóm phát triển nhỏ và lưu lượng ban đầu thấp, chia nhỏ dịch vụ làm tăng chi phí vận hành (nhiều tiến trình, mạng nội bộ, giao dịch phân tán) mà không mang lại lợi ích tương xứng. Ranh giới module của NestJS vẫn giữ được sự tách biệt về mã nguồn nếu sau này cần tách.

\begin{rtable}
\begin{longtable}{|L{3.6cm}|L{5.4cm}|L{5.8cm}|}
\caption{Các quyết định kiến trúc chính}\label{tab:arch-decisions}\\
\hline
\thead{Quyết định} & \thead{Lý do} & \thead{Đánh đổi} \\ \hline
\endfirsthead
\hline
\thead{Quyết định} & \thead{Lý do} & \thead{Đánh đổi} \\ \hline
\endhead
Một máy chủ NestJS phân module thay vì vi dịch vụ & Giao dịch đặt phòng -- thanh toán nằm trong một cơ sở dữ liệu; triển khai một container & Khó mở rộng độc lập từng phần; một lỗi nghiêm trọng ảnh hưởng toàn bộ API \\ \hline
MongoDB với giao dịch nhiều tài liệu & Lược đồ linh hoạt cho catalog; giao dịch ACID cho đặt phòng & Bắt buộc replica set; truy vấn liên kết (join) hạn chế \\ \hline
Chấm điểm LightGBM bằng TypeScript trong tiến trình & Không cần dịch vụ Python, không thêm độ trễ mạng & Phải bảo đảm khớp đặc trưng và điểm số giữa hai ngôn ngữ bằng kiểm thử \\ \hline
Gọi Open-Meteo trực tiếp từ iOS & Dữ liệu công khai, không cần khóa; giảm tải máy chủ & Không kiểm soát tập trung việc lưu đệm \\ \hline
Xác nhận thanh toán chỉ qua webhook & Không tin máy khách & Phụ thuộc vào việc payOS gọi được địa chỉ webhook công khai \\ \hline
JWT không trạng thái cho access token & Không truy vấn phiên cho mỗi yêu cầu & Thu hồi tức thì nhờ \code{tokenVersion} và \code{isActive}, vẫn cần đọc người dùng trong JwtStrategy \\ \hline
\end{longtable}
\end{rtable}

\subsection{Sơ đồ thành phần của máy chủ}

\Cref{fig:components} trình bày các module của máy chủ và quan hệ phụ thuộc giữa chúng. Module \code{Bookings} phụ thuộc \code{Hotels} để kiểm tra quyền sở hữu và phụ thuộc \code{Payments} để hủy liên kết thanh toán khi đặt phòng bị hủy. Module \code{Recommendations} không phụ thuộc module nghiệp vụ nào. Nó đọc trực tiếp các bộ sưu tập cần thiết qua kết nối MongoDB, nên lỗi trong phần gợi ý không ảnh hưởng tới luồng đặt phòng. Các thành phần dùng chung (guard, interceptor, filter, decorator) nằm trong thư mục \code{common} và được đăng ký toàn cục trong \code{AppModule}.

\begin{figure}[H]
\centering
\diagram{component-backend}
\caption{Sơ đồ thành phần của máy chủ NestJS}\label{fig:components}
\end{figure}

\subsection{Sơ đồ lớp}

\Cref{fig:class} là sơ đồ lớp rút gọn của các lớp dịch vụ chính, chỉ giữ các phương thức công khai có ý nghĩa nghiệp vụ. Các lớp controller không được vẽ vì chúng chỉ chuyển tiếp yêu cầu sang lớp dịch vụ tương ứng.

\begin{figure}[H]
\centering
\diagram{class-backend}
\caption{Sơ đồ lớp rút gọn của tầng dịch vụ}\label{fig:class}
\end{figure}

\subsection{Sơ đồ triển khai}\label{sec:deploy-design}

\Cref{fig:deployment} mô tả cách các thành phần được triển khai. Máy chủ API chạy trong một container Docker trên Render, phía trước có Cloudflare (phản hồi HTTP của \code{api.hotelia.asia} có tiêu đề \code{server: cloudflare} và \code{x-render-origin-server: Render}). Bảng điều khiển là ứng dụng tĩnh do Vercel phục vụ tại \code{admin.hotelia.asia}, có quy tắc chuyển mọi đường dẫn về \code{index.html} để định tuyến phía trình duyệt. Cơ sở dữ liệu là cụm MongoDB Atlas. Ảnh khách sạn được phục vụ bởi chính API tại đường dẫn \code{/media/}.

\begin{figure}[H]
\centering
\diagram{deployment}
\caption{Sơ đồ triển khai của Hotelia}\label{fig:deployment}
\end{figure}

\section{Thiết kế cơ sở dữ liệu}\label{sec:db-design}

\subsection{Mô hình thực thể -- liên kết}

\Cref{fig:erd} trình bày các bộ sưu tập (collection) và quan hệ giữa chúng. MongoDB không có khóa ngoại. Các quan hệ được biểu diễn bằng trường \code{ObjectId} tham chiếu và được kiểm tra ở tầng dịch vụ. Hai thiết kế đáng chú ý:
\begin{itemize}
\item Đặt phòng lưu \emph{bản sao} tên khách sạn, tên loại phòng và đơn giá tại thời điểm đặt (\code{hotelName}, \code{roomTypeName}, \code{roomPrice}). Nhờ vậy lịch sử đặt phòng không thay đổi khi chủ khách sạn đổi tên hay đổi giá sau này.
\item Tồn kho không được lưu như một số đếm giảm dần. Số phòng trống luôn được \emph{tính} từ tổng số phòng và các đặt phòng chồng lấn theo \cref{eq:free-rooms}. Cách này tránh sai lệch khi một đặt phòng hết hạn mà không có tác vụ nào cộng lại số phòng. Đổi lại, mỗi lần kiểm tra phải đọc các đặt phòng liên quan, nên chỉ mục \code{(roomTypeId, checkIn, checkOut)} là bắt buộc.
\end{itemize}

\begin{figure}[H]
\centering
\diagram{erd}
\caption{Sơ đồ thực thể -- liên kết của cơ sở dữ liệu Hotelia}\label{fig:erd}
\end{figure}

\subsection{Từ điển dữ liệu}

Các bảng dưới đây mô tả các trường của từng bộ sưu tập, được lấy từ lược đồ Mongoose và từ các thao tác ghi trực tiếp trong mã nguồn. Mọi bộ sưu tập dùng lược đồ Mongoose đều có thêm \code{createdAt} và \code{updatedAt} do tùy chọn \code{timestamps}.

\newcommand{\ddtable}[3]{%
\begin{rtable}
\begin{longtable}{|L{3.6cm}|L{2.4cm}|L{3.4cm}|L{5.4cm}|}
\caption{#1}\label{#2}\\
\hline
\thead{Trường} & \thead{Kiểu} & \thead{Ràng buộc} & \thead{Ý nghĩa} \\ \hline
\endfirsthead
\hline
\thead{Trường} & \thead{Kiểu} & \thead{Ràng buộc} & \thead{Ý nghĩa} \\ \hline
\endhead
#3
\end{longtable}
\end{rtable}}

\ddtable{Từ điển dữ liệu bộ sưu tập \code{users}}{tab:dd-users}{
\_id & ObjectId & khóa chính & Mã người dùng \\ \hline
email & String & bắt buộc, duy nhất, chữ thường & Email đăng nhập \\ \hline
passwordHash & String & bắt buộc & Băm bcrypt; tài khoản liên kết có giá trị ngẫu nhiên \\ \hline
googleSub, appleSub & String & duy nhất một phần & Định danh tại Google/Apple \\ \hline
name, phone, avatar & String & name bắt buộc & Hồ sơ \\ \hline
role & Enum & USER, OWNER, ADMIN & Vai trò \\ \hline
ownerStatus & Enum & PENDING, APPROVED, REJECTED & Trạng thái hồ sơ chủ khách sạn \\ \hline
businessName, businessLicense & String & & Thông tin doanh nghiệp khi đăng ký chủ khách sạn \\ \hline
isActive & Boolean & mặc định true & false khi bị quản trị viên khóa \\ \hline
loginAttempts, lockedUntil & Number, Date & & Chống dò mật khẩu \\ \hline
refreshToken, refreshTokenExpiry & String, Date & & Băm bcrypt của refresh token hiện hành \\ \hline
tokenVersion & Number & mặc định 0 & Tăng khi đổi/đặt lại mật khẩu để thu hồi token \\ \hline
passwordChangedAt, lastLoginAt & Date & & Mốc thời gian bảo mật \\ \hline
}

\ddtable{Từ điển dữ liệu bộ sưu tập \code{hotels}}{tab:dd-hotels}{
\_id & ObjectId & khóa chính & Mã khách sạn \\ \hline
ownerId & ObjectId & tham chiếu users, bắt buộc & Chủ khách sạn \\ \hline
name, description, address, city & String & bắt buộc & Thông tin cơ bản \\ \hline
district, country, latitude, longitude & String, Number & & Vị trí \\ \hline
status & Enum & DRAFT, PENDING\_APPROVAL, PUBLISHED, REJECTED, SUSPENDED & Trạng thái duyệt \\ \hline
starRating & Number & mặc định 0 (chưa rõ) & Hạng sao \\ \hline
averageRating, reviewCount & Number & mặc định 0 & Điểm từ đánh giá đã xác minh \\ \hline
demoRating, demoReviewCount, demoReviews & Number, Array & & Điểm và nhận xét biên tập (minh họa) \\ \hline
isDemoCatalog & Boolean & mặc định false & Khách sạn chỉ để xem, không cho đặt \\ \hline
amenities, images & [String] & & Tiện nghi, đường dẫn ảnh \\ \hline
checkInTime, checkOutTime & String & mặc định 14:00, 12:00 & Giờ nhận/trả phòng \\ \hline
rejectionReason & String & & Lý do từ chối \\ \hline
}

\ddtable{Từ điển dữ liệu bộ sưu tập \code{roomtypes}}{tab:dd-roomtypes}{
\_id & ObjectId & khóa chính & Mã loại phòng \\ \hline
hotelId & ObjectId & tham chiếu hotels, bắt buộc & Khách sạn sở hữu \\ \hline
name, description & String & bắt buộc & Tên và mô tả \\ \hline
basePrice & Number & bắt buộc, $> 0$ khi đặt & Giá một đêm (VND) \\ \hline
maxGuests & Number & bắt buộc & Sức chứa tối đa \\ \hline
totalRooms & Number & bắt buộc & Tồn kho (số phòng vật lý) \\ \hline
availableRooms & Number & & Giá trị hiển thị; khi đặt phòng luôn tính lại theo đêm \\ \hline
inventoryVersion & Number & & Tăng trong mỗi giao dịch đặt phòng để tuần tự hóa \\ \hline
amenities, images & [String] & & Tiện nghi phòng, ảnh \\ \hline
isActive & Boolean & mặc định true & false khi ngừng bán \\ \hline
}

\ddtable{Từ điển dữ liệu bộ sưu tập \code{bookings}}{tab:dd-bookings}{
\_id & ObjectId & khóa chính & Mã đặt phòng \\ \hline
userId, hotelId, roomTypeId & ObjectId & bắt buộc & Khách, khách sạn, loại phòng \\ \hline
checkIn, checkOut & Date & UTC 00:00 & Khoảng nửa mở $[d_{\text{in}}, d_{\text{out}})$ \\ \hline
nights, guestCount & Number & & Số đêm, số khách \\ \hline
guestName, guestEmail, guestPhone, specialRequests & String & & Thông tin liên hệ \\ \hline
roomPrice, totalPrice & Number & bắt buộc & Đơn giá và tổng tiền tại thời điểm đặt \\ \hline
hotelName, roomTypeName & String & & Bản sao để giữ lịch sử \\ \hline
status & Enum & 10 giá trị (\cref{tab:booking-states}) & Trạng thái \\ \hline
holdExpiresAt & Date & & Hạn giữ phòng (15 phút) \\ \hline
paymentMethod & Enum & PAY\_AT\_HOTEL, PAYOS & Phương thức đã chọn \\ \hline
requestKey, requestHash & String & duy nhất một phần với userId & Khóa lũy đẳng và băm nội dung \\ \hline
checkedInAt, checkedOutAt, cancelledAt & Date & & Mốc thời gian \\ \hline
cancelReason, refundAmount & String, Number & & Lý do hủy, số tiền hoàn (nếu có) \\ \hline
}

\ddtable{Từ điển dữ liệu bộ sưu tập \code{payments}}{tab:dd-payments}{
\_id & ObjectId & khóa chính & Mã thanh toán \\ \hline
bookingId, userId & ObjectId & bắt buộc & Đặt phòng, người trả \\ \hline
bookingKey & String & bắt buộc, duy nhất & Bảo đảm mỗi đặt phòng một thanh toán \\ \hline
amount & Number & bắt buộc & Số tiền lấy từ đặt phòng \\ \hline
method & Enum & PAYOS (giá trị khác chỉ để đọc dữ liệu cũ) & Kênh thanh toán \\ \hline
status & Enum & PENDING, PROCESSING, COMPLETED, FAILED, CANCELLED, REFUNDED & Trạng thái \\ \hline
orderCode & Number & duy nhất (sparse) & Mã đơn gửi payOS \\ \hline
transactionId, paymentUrl & String & & Mã liên kết, địa chỉ thanh toán \\ \hline
returnUrl, cancelUrl, expiresAt & String, Date & & Tham số liên kết \\ \hline
paidAt, failedAt, failureReason & Date, String & & Kết quả \\ \hline
metadata & String (JSON) & & Mã tham chiếu, thời điểm giao dịch từ webhook \\ \hline
}

\ddtable{Từ điển dữ liệu bộ sưu tập \code{reviews} và \code{notifications}}{tab:dd-reviews}{
\multicolumn{4}{|l|}{\textbf{reviews}} \\ \hline
hotelId, userId, bookingId & ObjectId & hotelId, userId bắt buộc & Khách sạn, người viết, đặt phòng làm căn cứ \\ \hline
rating & Number & 1--5 & Số sao \\ \hline
title, content & String & content bắt buộc & Nội dung \\ \hline
ownerReply & String & & Phản hồi của chủ khách sạn \\ \hline
helpfulCount, isVisible & Number, Boolean & isVisible mặc định true & Lượt hữu ích; false khi bị quản trị viên ẩn \\ \hline
\multicolumn{4}{|l|}{\textbf{notifications}} \\ \hline
userId & ObjectId & bắt buộc & Người nhận \\ \hline
title, message & String & bắt buộc & Nội dung \\ \hline
type, relatedId & String, ObjectId & & Loại (BOOKING, PAYMENT, \ldots), đối tượng liên quan \\ \hline
isRead & Boolean & mặc định false & Đã đọc \\ \hline
}

\ddtable{Từ điển dữ liệu các bộ sưu tập phụ trợ}{tab:dd-misc}{
\multicolumn{4}{|l|}{\textbf{favorites}: \_id = \code{userId:hotelId} (chống trùng)} \\ \hline
userId, hotelId, createdAt & ObjectId, Date & & Khách sạn yêu thích \\ \hline
\multicolumn{4}{|l|}{\textbf{support\_tickets}} \\ \hline
userId, subject, message & ObjectId, String & 3--160, 10--4000 ký tự & Yêu cầu hỗ trợ \\ \hline
status, reply, repliedAt & String, Date & OPEN, RESOLVED & Trạng thái, phản hồi của quản trị viên \\ \hline
\multicolumn{4}{|l|}{\textbf{password\_resets}: một tài liệu cho mỗi người dùng} \\ \hline
userId, tokenHash, expiresAt & ObjectId, String, Date & hạn 15 phút & SHA-256 của mã khôi phục, không lưu mã gốc \\ \hline
\multicolumn{4}{|l|}{\textbf{audit\_logs}} \\ \hline
actorId, role, method, path, createdAt & String, Date & & Ai thực hiện thao tác ghi nào; không lưu nội dung yêu cầu \\ \hline
\multicolumn{4}{|l|}{\textbf{recommendation\_views}: \_id = \code{actor:hotelId}} \\ \hline
actor, hotelId, viewedAt & String, ObjectId, Date & TTL 90 ngày & Lượt mở chi tiết, mỗi khách sạn một bản ghi cho mỗi tác nhân \\ \hline
\multicolumn{4}{|l|}{\textbf{recommendation\_impressions}: \_id = \code{actor:impressionId}} \\ \hline
candidateIds, candidatePositions & [ObjectId], [Object] & tối đa 8, không trùng & Thứ tự hiển thị \\ \hline
destinationKey, filters, searchTextPresent & String, Object, Boolean & & Ngữ cảnh truy vấn; chỉ giữ thành phố có trong danh mục, không lưu văn bản tìm kiếm tự do \\ \hline
detailOpenedHotelIds, detailOpenedAt & [ObjectId], Object & & Kết quả (khách sạn được mở) \\ \hline
rowLoadedAt, expiresAt & Date & TTL theo expiresAt (90 ngày) & Thời điểm hiển thị \\ \hline
\multicolumn{4}{|l|}{\textbf{roomavailabilities}: lược đồ được khai báo nhưng không còn được sử dụng} \\ \hline
}

\subsection{Chỉ mục}

\Cref{tab:indexes} liệt kê các chỉ mục được khai báo trong mã nguồn và truy vấn mà mỗi chỉ mục phục vụ.

\begin{rtable}
\begin{longtable}{|L{3.0cm}|L{5.4cm}|L{6.4cm}|}
\caption{Các chỉ mục của cơ sở dữ liệu}\label{tab:indexes}\\
\hline
\thead{Bộ sưu tập} & \thead{Chỉ mục} & \thead{Mục đích} \\ \hline
\endfirsthead
\hline
\thead{Bộ sưu tập} & \thead{Chỉ mục} & \thead{Mục đích} \\ \hline
\endhead
users & \code{\{role, ownerStatus\}}; \code{googleSub}, \code{appleSub} duy nhất một phần & Danh sách hồ sơ chờ duyệt; đăng nhập liên kết \\ \hline
hotels & \code{\{status, city\}}, \code{ownerId}, \code{averageRating}, văn bản trên name/description/city & Tìm kiếm khách sạn đã duyệt theo thành phố; khách sạn của chủ \\ \hline
roomtypes & \code{hotelId}, \code{basePrice} & Loại phòng của khách sạn; lọc theo giá \\ \hline
bookings & \code{\{roomTypeId, checkIn, checkOut\}} & Đếm đặt phòng chồng lấn (truy vấn nóng nhất) \\ \hline
bookings & \code{\{userId, requestKey\}} duy nhất một phần & Khóa lũy đẳng \\ \hline
bookings & \code{\{userId, createdAt\}}, \code{\{hotelId, checkIn\}}, \code{status} & Lịch sử của khách; đặt phòng của khách sạn \\ \hline
payments & \code{bookingKey} duy nhất; \code{orderCode} duy nhất sparse & Một thanh toán cho mỗi đặt phòng; tra cứu từ webhook \\ \hline
reviews & \code{\{hotelId, createdAt\}}, \code{userId}, \code{rating} & Đánh giá của khách sạn \\ \hline
notifications & \code{\{userId, isRead, createdAt\}} & Danh sách thông báo chưa đọc \\ \hline
recommendation\_* & \code{\{actor, viewedAt\}}, \code{\{actor, rowLoadedAt\}}, TTL & Lịch sử của tác nhân; tự xóa sau 90 ngày \\ \hline
\end{longtable}
\end{rtable}

\section{Thiết kế API}\label{sec:api-design}

\subsection{Quy ước chung}

API tuân theo phong cách REST với tiền tố \code{/api/v1}. Dữ liệu trao đổi là JSON. Yêu cầu cần xác thực mang tiêu đề \code{Authorization: Bearer <access token>}. Mọi lỗi được \code{HttpExceptionFilter} chuẩn hóa về cùng một cấu trúc:

\begin{lstlisting}[language=TypeScript, caption={Cấu trúc phản hồi lỗi thống nhất}, label={lst:error-format}, numbers=none]
{ "success": false,
  "error": { "code": "BAD_REQUEST", "message": "Loại phòng đã hết trong thời gian đã chọn",
             "statusCode": 400, "path": "/api/v1/bookings",
             "timestamp": "2026-10-24T14:57:52.291Z" } }
\end{lstlisting}

Thông điệp lỗi nghiệp vụ được viết bằng tiếng Việt ở máy chủ. Ứng dụng iOS dịch sang tiếng Anh khi khóa tương ứng có trong String Catalog. Nếu không có bản dịch, ứng dụng hiển thị một thông điệp chung bằng tiếng Anh thay vì để lẫn tiếng Việt (\code{APIError.errorDescription}).

\subsection{Danh mục điểm cuối}

\Cref{tab:api} liệt kê toàn bộ điểm cuối của API, được trích từ các controller. Cột ``Quyền'' dùng ký hiệu: \textbf{C} là công khai; \textbf{U} là mọi người dùng đã đăng nhập; \textbf{O} là chủ khách sạn; \textbf{A} là quản trị viên; \textbf{C$^{*}$} là công khai nhưng nhận token nếu có (gợi ý).

\begin{rtable}
\begin{longtable}{|C{1.7cm}|L{5.6cm}|C{1.3cm}|L{6.2cm}|}
\caption{Danh mục điểm cuối của API Hotelia}\label{tab:api}\\
\hline
\thead{Phương thức} & \thead{Đường dẫn} & \thead{Quyền} & \thead{Chức năng} \\ \hline
\endfirsthead
\hline
\thead{Phương thức} & \thead{Đường dẫn} & \thead{Quyền} & \thead{Chức năng} \\ \hline
\endhead
\multicolumn{4}{|l|}{\textbf{Xác thực -- \code{/auth}}} \\ \hline
POST & \code{/auth/register} & C & Đăng ký, trả cặp token \\ \hline
POST & \code{/auth/login} & C & Đăng nhập email -- mật khẩu \\ \hline
POST & \code{/auth/social} & C & Đăng nhập Google/Apple bằng ID token \\ \hline
POST & \code{/auth/refresh} & C & Xoay vòng refresh token \\ \hline
POST & \code{/auth/logout} & U & Thu hồi refresh token \\ \hline
POST & \code{/auth/forgot-password} & C & Gửi mã khôi phục qua dịch vụ email \\ \hline
POST & \code{/auth/reset-password} & C & Đặt mật khẩu mới bằng mã \\ \hline
\multicolumn{4}{|l|}{\textbf{Người dùng và tài khoản}} \\ \hline
GET, PUT & \code{/users/me} & U & Xem, cập nhật hồ sơ \\ \hline
GET & \code{/account/favorites} & U & Danh sách yêu thích \\ \hline
PUT, DELETE & \code{/account/favorites/:id} & U & Thêm, bỏ yêu thích \\ \hline
POST & \code{/account/password} & U & Đổi mật khẩu \\ \hline
GET, POST & \code{/account/support} & U & Xem, tạo yêu cầu hỗ trợ \\ \hline
GET & \code{/account/support/all} & A & Mọi yêu cầu hỗ trợ \\ \hline
POST & \code{/account/support/:id/reply} & A & Trả lời yêu cầu hỗ trợ \\ \hline
GET & \code{/notifications} & U & Thông báo (phân trang) \\ \hline
POST & \code{/notifications/:id/read, /notifications/read-all} & U & Đánh dấu đã đọc \\ \hline
\multicolumn{4}{|l|}{\textbf{Khách sạn và loại phòng}} \\ \hline
GET & \code{/hotels} & C & Tìm kiếm, lọc, sắp xếp, phân trang \\ \hline
GET & \code{/hotels/:id} & C & Chi tiết khách sạn đã duyệt \\ \hline
POST & \code{/hotels} & O, A & Tạo khách sạn (DRAFT) \\ \hline
GET & \code{/hotels/owner/my-hotels} & O, A & Khách sạn của tôi \\ \hline
PUT & \code{/hotels/:id} & O, A & Cập nhật khách sạn \\ \hline
POST & \code{/hotels/:id/submit} & O, A & Gửi duyệt \\ \hline
DELETE & \code{/hotels/:id} & O, A & Ngừng khách sạn (SUSPENDED) \\ \hline
GET & \code{/room-types/hotel/:hotelId?checkIn\&checkOut} & C & Loại phòng, kèm số phòng trống cho khoảng ngày \\ \hline
POST & \code{/room-types/hotel/:hotelId} & O, A & Tạo loại phòng \\ \hline
PUT, DELETE & \code{/room-types/:id} & O, A & Cập nhật (kiểm tra số phòng đã bán), ngừng bán \\ \hline
GET & \code{/room-types/:id/availability?start\&end} & U & Số phòng trống từng ngày (tối đa 90 ngày) \\ \hline
\multicolumn{4}{|l|}{\textbf{Đặt phòng và thanh toán}} \\ \hline
POST & \code{/bookings (tiêu đề Idempotency-Key)} & U & Tạo đặt phòng, giữ 15 phút \\ \hline
GET & \code{/bookings?status} & U & Đặt phòng của tôi \\ \hline
GET & \code{/bookings/:id} & U & Chi tiết (khách, chủ khách sạn hoặc quản trị viên) \\ \hline
POST & \code{/bookings/:id/pay-at-hotel} & U & Xác nhận trả tại khách sạn \\ \hline
POST & \code{/bookings/:id/cancel} & U & Hủy hoặc yêu cầu hủy \\ \hline
POST & \code{/bookings/:id/resolve-cancellation} & O, A & Xác nhận hủy \\ \hline
POST & \code{/bookings/:id/check-in, /check-out} & O & Nhận phòng, trả phòng \\ \hline
POST & \code{/payments} & U & Tạo hoặc lấy lại liên kết payOS \\ \hline
POST & \code{/payments/webhooks/payos} & C (ký số) & Webhook của payOS \\ \hline
GET & \code{/payments, /payments/:id, /payments/booking/:id} & U & Thanh toán của tôi \\ \hline
\multicolumn{4}{|l|}{\textbf{Đánh giá}} \\ \hline
GET & \code{/reviews/hotel/:hotelId} & C & Đánh giá đang hiển thị \\ \hline
POST & \code{/reviews} & U & Tạo đánh giá sau lưu trú \\ \hline
GET & \code{/reviews/my-reviews} & U & Đánh giá của tôi \\ \hline
POST & \code{/reviews/:id/helpful} & U & Đánh dấu hữu ích \\ \hline
GET & \code{/reviews/owner} & O, A & Đánh giá về khách sạn của tôi \\ \hline
POST & \code{/reviews/:id/reply} & O, A & Phản hồi đánh giá \\ \hline
POST & \code{/reviews/:id/hide} & A & Ẩn đánh giá, tính lại điểm \\ \hline
\multicolumn{4}{|l|}{\textbf{Chủ khách sạn -- \code{/owners}}} \\ \hline
POST & \code{/owners/apply} & U & Đăng ký làm chủ khách sạn \\ \hline
GET & \code{/owners/dashboard} & O, A & Thống kê tổng quan \\ \hline
GET & \code{/owners/bookings} & O, A & Mọi đặt phòng của các khách sạn của tôi \\ \hline
GET & \code{/owners/hotels/:hotelId/bookings} & O, A & Đặt phòng của một khách sạn \\ \hline
GET & \code{/owners/revenue?year} & O, A & Doanh thu theo tháng \\ \hline
\multicolumn{4}{|l|}{\textbf{Quản trị -- \code{/admin} (toàn bộ yêu cầu vai trò A)}} \\ \hline
GET & \code{/admin/stats} & A & Thống kê hệ thống \\ \hline
GET & /admin/users?page\&search\&role & A & Tìm kiếm người dùng \\ \hline
POST & \code{/admin/users/:id/suspend, /activate} & A & Khóa, mở khóa \\ \hline
GET & \code{/admin/owners, /admin/owners/pending} & A & Hồ sơ chủ khách sạn \\ \hline
POST & \code{/admin/owners/:id/approve, /reject} & A & Duyệt, từ chối \\ \hline
GET & \code{/admin/hotels, /admin/hotels/pending} & A & Khách sạn \\ \hline
POST & \code{/admin/hotels/:id/approve, /reject} & A & Duyệt, từ chối kèm lý do \\ \hline
GET & \code{/admin/bookings, /payments, /reviews, /audit-logs} & A & Giám sát (tối đa 200 bản ghi gần nhất) \\ \hline
\multicolumn{4}{|l|}{\textbf{Gợi ý -- \code{/recommendations}}} \\ \hline
GET & \code{/recommendations?sessionId\&destination\&}\ldots & C$^{*}$ & Danh sách gợi ý và thông tin xếp hạng \\ \hline
POST & \code{/recommendations/impressions} & C$^{*}$ & Ghi lượt hiển thị (tối đa 8 ứng viên) \\ \hline
POST & \code{/recommendations/views} & C$^{*}$ & Ghi lượt mở chi tiết \\ \hline
DELETE & \code{/recommendations/views?sessionId} & C$^{*}$ & Xóa lịch sử gợi ý của tác nhân \\ \hline
\multicolumn{4}{|l|}{\textbf{Vận hành}} \\ \hline
GET & \code{/health} & C & Kiểm tra sức khỏe (ngoài tiền tố /api) \\ \hline
\end{longtable}
\end{rtable}

\subsection{Tìm kiếm và lọc ứng viên}\label{sec:search-design}

Lớp \code{CatalogFilter} được hai điểm cuối \code{GET /hotels} và \code{GET /recommendations} dùng chung cho các tham số giá tối thiểu, giá tối đa, số khách, điểm tối thiểu, ngày nhận và ngày trả. Khi có bất kỳ tham số liên quan tới phòng, hàm \code{eligibleRooms} thực hiện ba bước: (i) tìm loại phòng đang bán, có sức chứa đủ và giá nằm trong khoảng; (ii) nếu có khoảng ngày, đọc các đặt phòng chồng lấn đang chiếm tồn kho; (iii) giữ lại loại phòng còn ít nhất một phòng ở mọi đêm. Một khách sạn được giữ lại khi còn ít nhất một loại phòng thỏa điều kiện.

Với điểm tối thiểu, khách sạn có đánh giá đã xác minh được so bằng \code{averageRating}; khách sạn chưa có đánh giá xác minh được so bằng điểm biên tập. Từ khóa điểm đến được thoát ký tự đặc biệt của biểu thức chính quy trước khi dùng, để tránh truy vấn có chi phí bất thường. Kết quả sắp theo điểm, hoặc theo giá thấp nhất khi người dùng chọn sắp xếp theo giá. Mỗi trang có tối đa 100 phần tử.

\section{Thiết kế phân quyền}\label{sec:rbac-design}

\Cref{tab:rbac} là ma trận phân quyền theo tài nguyên. Ký hiệu ``sở hữu'' nghĩa là thao tác chỉ thành công với tài nguyên thuộc người gọi. Điều kiện này được kiểm tra ở tầng dịch vụ, độc lập với kiểm tra vai trò ở guard.

\begin{rtable}
\begin{longtable}{|L{4.6cm}|C{2.2cm}|C{2.4cm}|C{2.6cm}|C{2.6cm}|}
\caption{Ma trận phân quyền theo vai trò}\label{tab:rbac}\\
\hline
\thead{Tài nguyên / thao tác} & \thead{Khách (chưa đăng nhập)} & \thead{USER} & \thead{OWNER} & \thead{ADMIN} \\ \hline
\endfirsthead
\hline
\thead{Tài nguyên / thao tác} & \thead{Khách (chưa đăng nhập)} & \thead{USER} & \thead{OWNER} & \thead{ADMIN} \\ \hline
\endhead
Xem khách sạn đã duyệt, loại phòng, đánh giá & \checkmark & \checkmark & \checkmark & \checkmark \\ \hline
Xem gợi ý & khám phá & cá nhân hóa & cá nhân hóa & cá nhân hóa \\ \hline
Tạo đặt phòng, thanh toán & -- & \checkmark & \checkmark & \checkmark \\ \hline
Xem, hủy đặt phòng & -- & sở hữu & sở hữu & xem tất cả \\ \hline
Đánh giá & -- & đã trả phòng & đã trả phòng & đã trả phòng \\ \hline
Tạo, sửa khách sạn, loại phòng & -- & -- & sở hữu & \checkmark \\ \hline
Nhận phòng, trả phòng & -- & -- & sở hữu & -- \\ \hline
Xác nhận hủy & -- & -- & sở hữu & \checkmark \\ \hline
Phản hồi đánh giá & -- & -- & sở hữu & \checkmark \\ \hline
Duyệt chủ khách sạn, khách sạn & -- & -- & -- & \checkmark \\ \hline
Khóa người dùng, ẩn đánh giá & -- & -- & -- & \checkmark \\ \hline
Thanh toán, nhật ký, thống kê toàn hệ thống & -- & -- & -- & \checkmark \\ \hline
\end{longtable}
\end{rtable}

\section{Thiết kế đường ống gợi ý}\label{sec:rec-design}

\subsection{Kiến trúc tổng quát}

Đường ống gợi ý của Hotelia gồm hai nửa tách biệt (\cref{fig:rec-pipeline}). Nửa \textbf{ngoại tuyến} chạy bằng Python trên máy phát triển: xuất ảnh chụp danh mục từ API, sinh dữ liệu huấn luyện, huấn luyện LightGBM LambdaRank, đánh giá và xuất mô hình ra JSON. Nửa \textbf{trực tuyến} chạy trong NestJS: lọc ứng viên, tổng hợp tín hiệu, tính 16 đặc trưng, chấm điểm bằng mô hình đã nạp và ghi nhận tương tác. Hai nửa giao tiếp qua đúng hai tệp: ảnh chụp danh mục (đầu vào huấn luyện) và tệp mô hình \code{hotelia-lgbm.active.json} (đầu ra huấn luyện, được commit cùng mã nguồn máy chủ). Tệp mô hình chứa thứ tự đặc trưng; máy chủ từ chối nạp nếu thứ tự này khác hằng \code{HOTELIA\_FEATURES} trong mã TypeScript.

\begin{figure}[H]
\centering
\diagram{rec-pipeline}
\caption{Kiến trúc đường ống gợi ý: huấn luyện ngoại tuyến và phục vụ trực tuyến}\label{fig:rec-pipeline}
\end{figure}

\subsection{Lọc ứng viên và tín hiệu người dùng}

Tập ứng viên là các khách sạn PUBLISHED có ít nhất một loại phòng đang bán, sau khi áp dụng điểm đến và bộ lọc phòng như \cref{sec:search-design}. Không giới hạn số ứng viên trước khi xếp hạng: mọi ứng viên hợp lệ đều được chấm điểm rồi mới cắt theo \code{limit} (mặc định 20). Ca kiểm thử TC-R03 kiểm chứng điều này.

Tín hiệu của một tác nhân gồm:
\begin{itemize}
\item \textbf{Lượt xem}: tối đa 50 lượt xem trong 90 ngày, mỗi khách sạn một bản ghi. Trọng số tính theo \cref{eq:decay}: $w = 2^{-d/14}$.
\item \textbf{Đặt phòng đã xác nhận}: tối đa 50 đặt phòng trong 90 ngày ở các trạng thái PAID, CONFIRMED, CHECKED\_IN, CHECKED\_OUT, COMPLETED. Mỗi khách sạn có trọng số 3, chỉ tính một lần.
\end{itemize}
Tác nhân là \code{user:<id>} khi có access token hợp lệ, ngược lại là \code{session:<uuid>} do ứng dụng sinh và lưu trên thiết bị. Khi người dùng đã đăng nhập, hệ thống không trộn lịch sử của phiên ẩn danh vào hồ sơ tài khoản (BR-22). Cách này tránh việc gán nhầm hành vi của người dùng chung thiết bị.

\subsection{Mười sáu đặc trưng}

\Cref{tab:features} định nghĩa 16 đặc trưng của mỗi cặp (truy vấn, khách sạn). Ký hiệu: $p$ là giá thấp nhất của khách sạn (thay bằng $\bar p$ nếu thiếu); $\bar p$ là giá trung bình của các ứng viên trong truy vấn; $\mathcal{S}$ là tập tín hiệu với trọng số $w_j$ và $W = \sum_j w_j$; $r$, $R$ lần lượt là số đánh giá và điểm trung bình; $\mathcal{A}$ là tập tiện nghi đã chuẩn hóa (bỏ dấu, chữ thường); $h$ là giờ hiện tại theo giờ Việt Nam.

\begin{rtable}
\begin{longtable}{|C{0.6cm}|L{3.3cm}|L{5.6cm}|L{5.2cm}|}
\caption{Mười sáu đặc trưng của mô hình xếp hạng}\label{tab:features}\\
\hline
\thead{\#} & \thead{Tên} & \thead{Công thức} & \thead{Ý nghĩa} \\ \hline
\endfirsthead
\hline
\thead{\#} & \thead{Tên} & \thead{Công thức} & \thead{Ý nghĩa} \\ \hline
\endhead
\multicolumn{4}{|l|}{\textit{Nhóm thuộc tính khách sạn}} \\ \hline
1 & min\_price\_millions & $p / 10^6$ & Mức giá tuyệt đối \\ \hline
2 & price\_relative & $p / \bar p$ & Đắt hay rẻ so với các lựa chọn khác trong cùng truy vấn \\ \hline
3 & star\_rating & $\max(0, \text{sao})$; 0 nếu chưa rõ & Hạng sao \\ \hline
4 & rating & $\max(0, R)$ & Điểm đánh giá \\ \hline
5 & review\_log & $\ln(1 + r)$ & Độ tin cậy của điểm, nén theo logarit \\ \hline
6 & amenity\_count & $|\mathcal{A}_{\text{raw}}|$ & Số tiện nghi khai báo \\ \hline
7 & quality\_prior & $\dfrac{r \cdot \min(5,R) + 5 \cdot 3{,}5}{(r + 5)\cdot 5}$ & Điểm chất lượng có prior Bayes: 5 đánh giá ảo ở mức 3,5 sao, tránh để một đánh giá 5 sao duy nhất đứng đầu \\ \hline
\multicolumn{4}{|l|}{\textit{Nhóm phù hợp với lịch sử (bằng 0 khi chưa có lịch sử)}} \\ \hline
8 & city\_affinity & $\frac{1}{W}\sum_j w_j \,[\text{city} = \text{city}_j]$ & Tỷ lệ trọng số lịch sử ở cùng thành phố \\ \hline
9 & price\_affinity & $\exp\!\big(-\big|\ln p - \frac{1}{W}\sum_j w_j \ln p_j\big|\big)$ & Độ gần với mức giá quen thuộc (theo tỷ lệ, không theo hiệu) \\ \hline
10 & star\_affinity & $\frac{1}{W}\sum_j w_j \max\!\big(0, 1 - \frac{|s - s_j|}{5}\big)$ & Độ gần về hạng sao \\ \hline
11 & amenity\_affinity & $\frac{1}{W}\sum_j w_j \frac{|\mathcal{A} \cap \mathcal{A}_j|}{|\mathcal{A} \cup \mathcal{A}_j|}$ & Độ tương đồng Jaccard về tiện nghi \\ \hline
12 & history\_count & $\min(50, |\mathcal{S}|)$ & Lượng lịch sử \\ \hline
13 & has\_history & $[W > 0]$ & Phân biệt khám phá và cá nhân hóa \\ \hline
\multicolumn{4}{|l|}{\textit{Nhóm ngữ cảnh thời gian (giờ Việt Nam)}} \\ \hline
14 & hour\_sin & $\sin(2\pi h / 24)$ & Mã hóa tuần hoàn giờ trong ngày: 23 giờ và 0 giờ gần nhau \\ \hline
15 & hour\_cos & $\cos(2\pi h / 24)$ & Như trên \\ \hline
16 & is\_weekend & $[\text{thứ Bảy hoặc Chủ nhật}]$ & Cuối tuần \\ \hline
\end{longtable}
\end{rtable}

\textbf{Ví dụ.} Một khách đã đặt một resort 5 sao ở Đà Nẵng giá 1,2 triệu đồng ($w = 3$), đang xem danh sách ở Đà Nẵng với giá trung bình 1,0 triệu. Một ứng viên 4 sao ở Đà Nẵng giá 1,1 triệu có \code{price\_relative} $= 1{,}1$, \code{city\_affinity} $= 1$, \code{price\_affinity} $= e^{-|\ln 1{,}1 - \ln 1{,}2|} \approx 0{,}917$ và \code{star\_affinity} $= 1 - 1/5 = 0{,}8$. Một ứng viên 3 sao ở Hội An giá 0,6 triệu có \code{city\_affinity} $= 0$, \code{price\_affinity} $= 0{,}5$ và \code{star\_affinity} $= 0{,}6$.

\subsection{Dữ liệu huấn luyện trước phát hành}\label{sec:training-data}

Vì chưa có đủ hành vi người dùng thật, dữ liệu huấn luyện được sinh có kiểm soát. Phiên bản v3 sinh dữ liệu như sau:
\begin{enumerate}
\item \textbf{Danh mục}: dùng nguyên 100 khách sạn trong ảnh chụp API, với giá, hạng sao, điểm, số đánh giá và tiện nghi đúng như lúc phục vụ.
\item \textbf{Người dùng mô phỏng}: chọn ngẫu nhiên một thành phố và một khách sạn ``ưa thích''. 76\% truy vấn có lịch sử. Số tín hiệu là 1--12 (70\%) hoặc 13--50 (30\%). Khách sạn trong lịch sử được chọn theo phân phối softmax của độ hợp về thành phố, giá, sao và tiện nghi. 18\% tín hiệu là đặt phòng (trọng số 3), còn lại là lượt xem có tuổi ngẫu nhiên 0--90 ngày.
\item \textbf{Nhóm ứng viên}: 45\% truy vấn có 8 ứng viên (4 cùng thành phố, 4 ngẫu nhiên); 35\% là toàn bộ khách sạn của một thành phố (trung bình 8,8); 20\% là cả 100 khách sạn, tương ứng truy vấn không chọn điểm đến.
\item \textbf{Nhãn}: lợi ích kỳ vọng của mỗi ứng viên là tổ hợp tuyến tính
$U = 1{,}75\,x_8 + 1{,}0\,x_9 + 0{,}65\,x_{10} + 0{,}55\,x_{11} + 0{,}7\,x_7 - 0{,}12\,x_2$, cộng 0,12 cho tiện nghi gia đình vào cuối tuần. Cộng nhiễu Gumbel với tham số tỷ lệ 0,34 rồi chọn ứng viên lớn nhất. Cách làm này tương đương lựa chọn theo mô hình logit đa thức. Ứng viên được chọn có nhãn 1, các ứng viên khác có nhãn 0.
\end{enumerate}
Các hệ số trên là \emph{giả định} của người thiết kế dữ liệu, không được ước lượng từ người dùng thật. \Cref{sec:eval-model} phân tích hệ quả của giả định này đối với kết quả đánh giá.

\subsection{Mô hình và tiêu chí chọn}

Mô hình là \code{LGBMRanker} với mục tiêu \code{lambdarank}, tối đa 1000 cây, tốc độ học 0,05, 31 lá mỗi cây, tối thiểu 50 mẫu mỗi lá, lấy mẫu 90\% đặc trưng và 90\% dữ liệu mỗi vòng. Dữ liệu chia theo truy vấn thành 70\% huấn luyện, 15\% kiểm định và 15\% kiểm tra. Số cây được chọn bằng dừng sớm theo NDCG@10 trên tập kiểm định (dừng sau 50 vòng không cải thiện). Mọi chỉ số báo cáo đều tính trên tập kiểm tra, vốn không tham gia vào bất kỳ quyết định huấn luyện nào.

\subsection{Cơ chế dự phòng và ghi nhận tương tác}

Bộ xếp hạng dự phòng trong \code{ranking.ts} cho điểm
\begin{equation}
  \text{score} = \begin{cases}
    0{,}7 \cdot \frac{1}{W}\sum_j w_j\, \text{sim}(h, h_j) + 0{,}3 \cdot \text{quality\_prior} & \text{nếu } W > 0,\\
    \text{quality\_prior} & \text{nếu } W = 0,
  \end{cases}
  \label{eq:content-score}
\end{equation}
với $\text{sim} = 0{,}4\,[\text{cùng thành phố}] + 0{,}25\,\text{Jaccard tiện nghi} + 0{,}2\,\frac{\min(p, p_j)}{\max(p, p_j)} + 0{,}15\,(1 - |s - s_j|/5)$. Sau đó các khách sạn được chọn tham lam, mỗi lần trừ điểm theo số khách sạn cùng thành phố đã được chọn: 0,08 cho chế độ khám phá và 0,015 khi có lịch sử. Bước này tạo sự đa dạng thành phố cho người dùng mới, tương tự ý tưởng xếp hạng lại theo độ đa dạng \cite{carbonell1998}. Bộ xếp hạng dự phòng được dùng khi tệp mô hình không tồn tại, sai thứ tự đặc trưng, không có cây, hoặc khi chấm điểm gặp lỗi. Phản hồi của API cho biết bộ xếp hạng nào đã được dùng (\code{ranking.modelUsed}, \code{ranking.version}).

Ứng dụng ghi nhận hai loại sự kiện. \textbf{Lượt hiển thị} (impression) là thứ tự tối đa 8 khách sạn đầu tiên của hàng gợi ý, kèm ngữ cảnh truy vấn. \textbf{Lượt mở chi tiết} chỉ được gắn vào một lượt hiển thị khi đúng tác nhân, đúng mã lượt hiển thị và khách sạn có trong danh sách. Thiết kế này nhằm sau này tạo được nhãn ``hiển thị nhưng không mở'' và ``hiển thị và mở'' từ hành vi thật. Để bảo vệ quyền riêng tư, hệ thống không lưu văn bản tìm kiếm tự do (chỉ lưu thành phố nếu khớp danh mục), dữ liệu tự xóa sau 90 ngày, và người dùng có thể tự xóa lịch sử (UC12). Các thiết kế này nhằm đáp ứng nguyên tắc tối thiểu hóa dữ liệu của Nghị định 13/2023/NĐ-CP \cite{nd13}. Cơ chế xin đồng ý trước khi thu thập chưa được cài đặt và được nêu trong hướng phát triển.

\section{Thiết kế giao diện người dùng song ngữ}\label{sec:ui-design}

\subsection{Nguyên tắc thiết kế}

Giao diện iOS tuân theo hướng dẫn giao diện người của Apple \cite{applehig}: thanh tab dưới cùng cho ba khu vực chính, bảng (sheet) trượt từ dưới lên cho các tác vụ ngắn như đặt phòng hay chọn ngôn ngữ, và nút hành động chính đặt ở cuối màn hình. Bảng điều khiển web dùng bố cục thanh bên cố định, phân nhóm ``Vận hành'' và ``Theo dõi'' cho chủ khách sạn, ``Điều hành'' và ``Kiểm soát'' cho quản trị viên. Ba nguyên tắc nội dung được áp dụng nhất quán:
\begin{itemize}
\item \textbf{Trung thực về dữ liệu}: điểm biên tập luôn kèm nhãn ``Điểm Hotelia'' và ghi chú không phải đánh giá của khách đã lưu trú; khách sạn chỉ để xem không có nút đặt phòng; ảnh minh họa điểm đến được gắn nhãn ``Ảnh điểm đến''.
\item \textbf{Nói rõ điều sẽ xảy ra}: màn đặt phòng ghi rõ thời gian giữ 15 phút, giá cuối được xác nhận khi tạo đặt phòng, và ``chưa có cam kết hoàn tiền tự động''.
\item \textbf{Trạng thái nhìn thấy được}: mỗi đặt phòng có nhãn trạng thái màu (chờ thanh toán màu cam, đã xác nhận màu xanh dương, đã hủy màu đỏ).
\end{itemize}

\subsection{Cấu trúc điều hướng}

\Cref{fig:nav-ios} trình bày sơ đồ điều hướng của ứng dụng iOS và \cref{fig:nav-web} trình bày sơ đồ trang của bảng điều khiển. Bảng điều khiển định tuyến theo vai trò: \code{/owner/*} cho OWNER và ADMIN, \code{/admin/*} chỉ cho ADMIN. Người dùng không đủ quyền được chuyển về trang đăng nhập.

\begin{figure}[H]
\centering
\diagram{nav-ios}
\caption{Sơ đồ điều hướng của ứng dụng iOS}\label{fig:nav-ios}
\end{figure}

\begin{figure}[H]
\centering
\diagram{nav-web}
\caption{Sơ đồ trang của bảng điều khiển chủ khách sạn và quản trị viên}\label{fig:nav-web}
\end{figure}

\subsection{Thiết kế song ngữ}

Ứng dụng iOS có ba lựa chọn ngôn ngữ: theo hệ thống, tiếng Việt và tiếng Anh. Lựa chọn được lưu trong \code{UserDefaults} (khóa \code{appLanguage}). Hàm \code{L10n.text} đọc chuỗi từ gói ngôn ngữ tương ứng trong String Catalog, hiện có 920 khóa. Ngày, số và tiền tệ được định dạng theo \code{Locale} của ngôn ngữ đã chọn. Ví dụ, cùng một số tiền hiển thị là ``520.933'' trong tiếng Việt và ``520,933'' trong tiếng Anh. Khi người dùng đổi ngôn ngữ, toàn bộ cây giao diện được dựng lại để không còn chuỗi cũ. Thông điệp lỗi do máy chủ trả về được dịch ở phía ứng dụng nếu có khóa tương ứng.

Hai giới hạn được ghi nhận. Thứ nhất, nội dung dữ liệu như mô tả khách sạn và nhận xét biên tập phần lớn chỉ có tiếng Việt, nên vẫn hiển thị tiếng Việt khi giao diện ở tiếng Anh. Thứ hai, bảng điều khiển web hiện chỉ có tiếng Việt; bản tiếng Anh là hướng phát triển.

\subsection{Hệ thống thiết kế}

Ứng dụng iOS định nghĩa màu, khoảng cách và kiểu chữ dưới dạng hằng số dùng chung (\code{AppColors}, \code{AppSpacing}, \code{AppTypography}). Màu chính là xanh dương \code{\#007AFF}; các màu trạng thái gồm xanh lá \code{\#34C759}, đỏ \code{\#FF3B30} và cam \code{\#FF9500}. Bảng điều khiển định nghĩa token thiết kế trong \code{styles/tokens.css} và dùng Tailwind CSS. Hai nền tảng có giao diện riêng, nhưng cùng nguyên tắc nhãn trạng thái và cùng thuật ngữ nghiệp vụ.

\section{Kết luận chương}

Chương này đã phân tích yêu cầu của Hotelia với 37 yêu cầu chức năng, 16 yêu cầu phi chức năng và 37 ca sử dụng. Chương cũng đặc tả 25 quy tắc nghiệp vụ của luồng đặt phòng, máy trạng thái của đặt phòng, khách sạn và thanh toán, cùng các biểu đồ tuần tự chính. Về cấu trúc, chương trình bày kiến trúc phân lớp, cơ sở dữ liệu 13 bộ sưu tập, 78 điểm cuối API, ma trận phân quyền, đường ống gợi ý với 16 đặc trưng và thiết kế giao diện song ngữ. \Cref{chap:implementation} trình bày cách các thiết kế này được cài đặt.

